\begingroup
\small
\begin{thebibliography}{99}
\setlength{\itemsep}{0pt}
\setlength{\parskip}{0pt}
\addcontentsline{toc}{chapter}{Références}
\bibitem{sujet} \emph{Projet d'évaluation : Un éditeur de documents partagé. Développement web et cloud.} Sujet fourni par le porteur du projet, 3 pages. Fichier local : \texttt{sujet\_projet\_stage-1.pdf}.
\bibitem{cours} \emph{Ingénierie des exigences. Lesson 2 -- Spécification.} Support fourni, 39 diapositives, fondé sur Ian Sommerville, \emph{Software Engineering}, 9e édition. Fichier source : \texttt{Lecon-2\_GL-Spécification.pdf}. Les numéros cités désignent les diapositives.
\bibitem{cours-aws} \emph{Introduction à AWS.} Support de cours fourni, 117 diapositives. Fichier local \texttt{introduction-aws.pdf}, modules sur IAM, EC2, groupes de sécurité, S3, RDS, ALB et architectures à deux zones ; consulté le 26 octobre 2026.
\bibitem{decisions-porteur} Porteur du projet. \emph{Cadrage de l'historique, des versions métier, des branches, des droits, du gel quotidien et de la vérification e-mail.} Décisions communiquées le 25 octobre 2026. Ces précisions et extensions ne sont pas attribuées au sujet initial.
\bibitem{google-ot} John Day-Richter, Google. \emph{What's different about the new Google Docs: Conflict resolution.} 22 septembre 2010. Source historique.\par\url{https://drive.googleblog.com/2010/09/whats-different-about-new-google-docs_22.html}
\bibitem{google-protocole} John Day-Richter, Google. \emph{What's different about the new Google Docs: Making collaboration fast.} 23 septembre 2010. Source historique.\par\url{https://drive.googleblog.com/2010/09/whats-different-about-new-google-docs.html}
\bibitem{google-partage} Google, Aide Éditeurs Google Docs. \emph{Partager des fichiers via Google Drive.} Fonctions et limites de partage.\par\url{https://support.google.com/docs/answer/2494822?hl=fr}
\bibitem{google-historique} Google, Aide Éditeurs Google Docs. \emph{Identifier les modifications apportées à un fichier.}\par\url{https://support.google.com/docs/answer/190843?hl=fr}
\bibitem{google-export} Google, Aide Éditeurs Google Docs. \emph{Créer, afficher ou télécharger un fichier.}\par\url{https://support.google.com/docs/answer/49114?hl=fr}
\bibitem{overleaf-collaboration} Overleaf. \emph{Collaborating in Overleaf.} Documentation du mécanisme OT et de WebSocket.\par\url{https://docs.overleaf.com/collaborating/collaborating-in-overleaf}
\bibitem{overleaf-code} Overleaf. Dépôt officiel \texttt{overleaf/overleaf}, manifeste racine \texttt{package.json}, branche \texttt{main}, état consulté le 24 octobre 2026. Identifiant Git du blob observé : \texttt{8965f9ccbd7d772f1f32c33e1e013406ba238473}. Les dépendances peuvent évoluer.\par\url{https://github.com/overleaf/overleaf/blob/main/package.json}
\bibitem{overleaf-microservices} Overleaf. \emph{Microservices.} Documentation des éditions sur site, section architecture et processus d'édition/compilation.\par\url{https://docs.overleaf.com/on-premises/getting-started/microservices}
\bibitem{overleaf-infra} Overleaf. \emph{Server Pro infrastructure.} Documentation de déploiement sur site.\par\url{https://docs.overleaf.com/on-premises/getting-started/server-pro-infrastructure}
\bibitem{overleaf-backups} Overleaf. \emph{Data and backups.} Documentation de persistance et sauvegarde sur site.\par\url{https://docs.overleaf.com/on-premises/maintenance/data-and-backups}
\bibitem{yjs-intro} Yjs. \emph{Introduction.} Documentation officielle du modèle CRDT et de ses intégrations.\par\url{https://docs.yjs.dev/}
\bibitem{yjs-updates} Yjs. \emph{Document Updates.} Documentation officielle des propriétés et de l'API des mises à jour.\par\url{https://docs.yjs.dev/api/document-updates}
\bibitem{etherpad-release} Etherpad. \emph{Version 3.3.6.} Publication officielle, 21 septembre 2026.\par\url{https://github.com/ether/etherpad/releases/tag/v3.3.6}
\bibitem{etherpad-scale} Etherpad. \emph{Scale Calculator.} Calculateur de ressources, qui précise que ses résultats sont des repères approximatifs et que certaines performances de base de données distribuée n'ont pas été testées. Consulté le 26 octobre 2026.\par\url{https://scale.etherpad.org/}
\bibitem{etherpad-loadtest} Etherpad. \emph{etherpad-load-test.} Dépôt officiel de l'outil de simulation de charge ; les consignes distinguent auteurs actifs et lecteurs et recommandent d'adapter le scénario au cas d'usage. Consulté le 26 octobre 2026.\par\url{https://github.com/ether/etherpad-load-test}
\bibitem{etherpad-7756} Etherpad. \emph{Issue \#7756, Performance: Look deeper into scaling N+ contributors per pad.} Ouverte le 15 mai 2026. Discussion de pistes de recherche, sans résultats du test PostgreSQL/PgBouncer transmis avec la demande. Consultée le 26 octobre 2026.\par\url{https://github.com/ether/etherpad/issues/7756}
\bibitem{nodejs20-cli} Node.js. \emph{Documentation v20.20.2, Command-line API.} Description de \texttt{UV\_THREADPOOL\_SIZE} et des API qui utilisent le pool libuv ; consultée le 26 octobre 2026.\par\url{https://nodejs.org/docs/latest-v20.x/api/cli.html\#uv\_threadpool\_sizesize}
\bibitem{better-auth-otp} Better Auth. \emph{Email OTP.} Documentation officielle du plugin intégré pour la connexion par code et la vérification de l'adresse e-mail, consultée le 25 octobre 2026.\par\url{https://better-auth.com/docs/plugins/email-otp}
\bibitem{better-auth-fastify} Better Auth. \emph{Fastify Integration Guide.} Handler et lecture des sessions, consulté le 25 octobre 2026.\par\url{https://better-auth.com/docs/integrations/fastify}
\bibitem{better-auth-postgres} Better Auth. \emph{PostgreSQL Adapter.} Configuration de la base et schéma dédié, consulté le 25 octobre 2026.\par\url{https://better-auth.com/docs/adapters/postgresql}
\bibitem{postgres-doc} PostgreSQL Global Development Group. \emph{Documentation PostgreSQL 15.} Transactions et contraintes.\par\url{https://www.postgresql.org/docs/15/}
\bibitem{nodemailer-doc} Nodemailer. \emph{Nodemailer documentation.} Transport SMTP configurable.\par\url{https://nodemailer.com/smtp}
\bibitem{fastify} Fastify. \emph{Migration Guide V5} et publication v5.12.7. Node.js 20 ou supérieur est pris en charge par la branche 5.\par\url{https://fastify.dev/docs/v5.10.x/Guides/Migration-Guide-V5/}\par\url{https://github.com/fastify/fastify/releases/tag/v5.12.7}
\bibitem{aws-cost-explorer} Amazon Web Services. \emph{Exploring your data using Cost Explorer.} Documentation de consultation des coûts, consultée le 26 octobre 2026.\par\url{https://docs.aws.amazon.com/cost-management/latest/userguide/ce-exploring-data.html}
\bibitem{aws-local-artifacts} Porteur du projet. Fichiers consultés le 26 octobre 2026 : \texttt{application/infra/aws.yaml}, \texttt{application/infra/README.md}, \texttt{application/ci-local.sh}, \texttt{.github/workflows/ci.yml}, \texttt{.github/workflows/cd.yml}, \texttt{application/backend/src/server.ts}. Sources de l'état configuré et des routes ; les essais AWS n'ont pas tous été rejoués pour cette édition.
\bibitem{aws-test-records} Porteur du projet. \emph{CONTEXTE\_CLAUDE\_LOCAL.md}, sections « Premier déploiement AWS », « Point d'étape » et « Tests réels sur AWS ». Relevé local non versionné, consulté le 26 octobre 2026. Il conserve les résultats rapportés ; le fichier n'est pas inclus dans les livrables du dépôt.
\end{thebibliography}

\endgroup
